%\documentclass[12pt]{scrreprt}
\documentclass[12pt]{report}

% language may be romanian or english (default is english)
% type may be bachelor or master (default is bachelor)
\usepackage[language=english, type=bachelor]{style}

%controlling the appearance of your headers and footers
\usepackage{fancyhdr}
\pagestyle{fancy}
\lhead{}
\chead{}
\renewcommand{\headrulewidth}{0.2pt}
\renewcommand{\footrulewidth}{0.2pt}

\usepackage{listings}
\usepackage{xcolor}

\lstset{
  basicstyle=\ttfamily\small,
  keywordstyle=\color{blue}\bfseries,
  commentstyle=\color{gray},
  stringstyle=\color{red!70!black},
  breaklines=true,
  frame=single,
  captionpos=b,
  tabsize=2
}

\begin{document}

\specialization{Computer Science}
\title{Design and Evaluation of a Trustworthy AI Shopping Assistant
for E-Commerce}
\author{Mihai}
\supervisor{[Title and name of supervisor]}

\maketitle

\newpage
\thispagestyle{empty}
\mbox{}
\newpage
\pagenumbering{roman}

\cleardoublepage
\textbf{\Large ABSTRACT}
\vspace{0.5cm}
\hrule
\vspace{0.5cm}

Online shopping systems increasingly integrate large language models
to support product discovery and decision-making. However, many AI
shopping assistants are intrusive, poorly timed, or perceived as
upselling tools rather than genuine advisors. This thesis presents
the design, implementation, and evaluation of a reactive-first AI
shopping advisor for e-commerce. The proposed system combines a
classical recommendation layer based on co-occurrence, content
similarity, and personalized ranking with a behavior-aware readiness
model and an LLM-based explanation layer. Instead of proactively
interrupting the user, the system uses lightweight browsing signals
to infer when help may be useful and only generates AI advice when
the user explicitly engages.

The system is implemented as a full-stack prototype with a simulated
e-commerce store, a seeded product and order dataset, a
recommendation backend, and an interactive advisor interface.
Evaluation is performed at three levels: recommendation quality,
system performance, and user perception. Offline recommender metrics
are used to assess the baseline recommendation layer, while latency
and token-based cost measurements capture the efficiency of the
AI-assisted pipeline. A structured user evaluation compares the
baseline shopping experience with the reactive-first AI-assisted
version in terms of perceived usefulness, trust, intrusiveness, and
purchase confidence.

The thesis aims to show that hybrid architectures can deliver useful
and explainable shopping assistance while avoiding the drawbacks of
overly proactive AI interaction. The results contribute practical
guidance for building trustworthy, non-intrusive AI assistants in
consumer-facing systems.

\tableofcontents

\newpage
\pagenumbering{arabic}

\chapter{Introduction}
\label{chap:intro}

\section{Problem Context}
\label{sec:intro-context}

\par E-commerce has become one of the most important branches of the digital economy, with accelerated growth over the past decade. E-commerce platforms provide consumers with access to millions of products, but this abundance creates a fundamental problem: decision overload. Users spend significant time browsing, comparing, and evaluating products before making a purchase decision, and this process becomes increasingly complex as the supply grows.

\par Recommendation systems have emerged as a natural solution to this problem, employing various algorithmic techniques --- from collaborative filtering and co-occurrence analysis to content-based similarity --- to reduce the search space and present users with relevant products \cite{LLM4Rec2025, ComprehensiveRecSys2024}. These classical systems are computationally efficient and well-studied in the literature, but they present an important limitation: their outputs are opaque. A user who receives a ``similar products'' or ``frequently bought together'' recommendation does not understand why those products are relevant to their specific situation.

\par With the emergence of large language models (LLMs), it has become possible to generate natural language explanations for recommendations, to synthesize other buyers' reviews, and to provide structured comparisons between similar products \cite{LLMEnhancedRecSys2024, ExplainableRecKG2024}. However, integrating LLMs into commercial applications raises significant challenges related to trust, accuracy, and interaction design.

\par A relevant example is Amazon Rufus, Amazon's AI shopping assistant, which serves over 300 million users and has generated estimated incremental sales of \$12 billion \cite{RufusRevenue2025}. Despite these impressive figures, independent studies have shown that Rufus has an accuracy of only 32\% in identifying the most suitable products, a price hallucination rate of 28\%, and an 83\% tendency to recommend Amazon-owned products \cite{RufusWrong2024, RufusConfidentlyWrong2024}. These documented problems underscore the fact that, even at the largest scale, challenges related to factual accuracy, user trust, and transparency remain unsolved.

\par At the consumer perception level, the data is equally concerning: 55\% of consumers state that they do not trust AI chatbot recommendations \cite{ConsumerDistrust2025}, and 42\% perceive e-commerce AI assistants as upselling tools rather than genuine advisors \cite{AIUpsell2025}. Research in human-computer interaction (HCI) suggests that AI interventions timed after a natural user pause are more effective than proactive ones \cite{CHI2024Timing}, which opens the path for reactive interaction models.


\section{Motivation}
\label{sec:intro-motivation}

\par The combination of classical recommendation systems, which excel at fast candidate generation, and LLM models, which can generate explanations and summaries, offers an opportunity to build assistance systems that are simultaneously relevant, explainable, and efficient. However, the practical integration of these components raises essential questions: when should the system intervene? How can it justify recommendations without introducing fabricated information? What architecture provides the best tradeoff between relevance, latency, cost, and user trust?

\par This thesis is motivated by the observation that existing AI shopping assistants tend to be either too intrusive --- intervening automatically without the user's explicit consent --- or insufficiently transparent in how they generate recommendations. This work proposes a \textit{reactive-first} interaction model, in which the system monitors behavioral signals from the user (time spent on a product, scroll depth, interaction with the reviews section) to detect moments of readiness for assistance, but generates AI advice only when the user explicitly requests it.

\par This approach is supported by recent research showing that AI suggestions offered after a natural pause are better received than preemptive ones \cite{CHI2024Timing}, as well as by market data indicating significant consumer resistance to proactive AI chatbots \cite{ConsumerResistanceAI}.


\section{Objectives}
\label{sec:intro-objectives}

\par The main objective of this thesis is the design, implementation, and evaluation of a trustworthy AI shopping assistant that combines classical recommendation methods with explanations generated by large language models in a reactive interaction model. Specifically, the thesis aims to:

\begin{enumerate}
  \item \textbf{Design a hybrid recommendation architecture} that uses classical algorithms (co-occurrence with Jaccard coefficient, content-based similarity, trends with exponential time decay) for fast candidate generation, and an LLM model (Gemini Flash) for re-ranking and generating natural language explanations.

  \item \textbf{Implement a reactive interaction model} based on lightweight behavioral signals --- dwell time, scroll depth, review interaction --- that indicate the user's readiness for assistance without automatically triggering LLM calls.

  \item \textbf{Develop a review summarization mechanism} grounded in real buyer data, which surfaces both majority sentiments and conflicting opinions, addressing the most criticized problem of existing systems such as Rufus.

  \item \textbf{Implement an AI-assisted comparison mode} that allows users to compare 2--3 similar products in a structured format with tradeoff analysis.

  \item \textbf{Evaluate the system} at three levels: recommendation quality (offline metrics), system performance (latency and cost), and user perception (usefulness, trust, intrusiveness) through a structured user study.
\end{enumerate}


\section{Research Questions}
\label{sec:intro-rq}

\par To guide the design and evaluation of the system, this thesis formulates the following research questions:

\begin{description}
  \item[\textbf{RQ1:}] Does a reactive AI assistant produce a higher level of trust and a lower perception of intrusiveness compared to proactive AI suggestions?

  \item[\textbf{RQ2:}] Do LLM summaries grounded in buyer reviews improve the perceived usefulness and credibility of product information?

  \item[\textbf{RQ3:}] Does AI-assisted comparison improve user confidence when choosing between similar products?

  \item[\textbf{RQ4:}] What are the latency and cost tradeoffs of a two-stage recommendation architecture?
\end{description}


\section{Contributions}
\label{sec:intro-contributions}

\par The main contributions of this thesis are:

\begin{enumerate}
  \item A working \textbf{hybrid recommendation architecture} that combines classical algorithms with an LLM explanation layer in a two-stage pipeline (candidate generation under 50ms, LLM re-ranking and explanation under 500ms).

  \item A \textbf{reactive-first interaction model} based on behavioral signals, where the system prepares context and displays subtle indicators, and the LLM call is triggered only at the user's initiative.

  \item A \textbf{review summarization mechanism} grounded in real data, which surfaces conflicts and divergent opinions, not just the majority sentiment.

  \item An \textbf{AI-assisted comparison flow} for similar products within the same category.

  \item A \textbf{complete working prototype} implemented as a full-stack web application with a simulated store, realistic seed data, a recommendation engine, behavioral tracking, and an advisor interface.

  \item A \textbf{systematic evaluation} at three levels: offline recommendation metrics, performance measurements (latency, token cost), and a structured user study.

  \item \textbf{Design guidelines} for building non-intrusive AI assistants in consumer-facing systems.
\end{enumerate}


\section{Thesis Structure}
\label{sec:intro-structure}

\par The thesis is organized into seven chapters, as follows:

\par \textbf{Chapter~\ref{chap:background} --- Background and Related Work} presents the theoretical foundations of classical recommendation systems, the integration of LLMs into recommendation systems, research on trust and explainability in human-AI interaction, and a detailed case study of Amazon Rufus as a reference point for production-scale AI shopping assistants.

\par \textbf{Chapter~\ref{chap:problem} --- Problem Definition and Requirements} formulates the concrete problem addressed by this thesis, defines the functional and non-functional system requirements, establishes the evaluation criteria, and clarifies the project's boundaries and assumptions.

\par \textbf{Chapter~\ref{chap:design} --- System Design} describes the overall system architecture, the two-stage recommendation pipeline, the reactive interaction model, the behavioral tracking mechanism, the review summarization flow and comparison mode, as well as the justification for architectural decisions.

\par \textbf{Chapter~\ref{chap:implementation} --- Implementation} presents the technical details of the prototype: the technology stack, the monorepo structure, the Convex data model, the implemented recommendation algorithms, the LLM integration with Gemini Flash, and the user interface.

\par \textbf{Chapter~\ref{chap:evaluation} --- Evaluation} describes the evaluation methodology and presents results at all three levels: offline recommendation metrics, system performance measurements, and user study results. It includes a discussion of the results, limitations, and threats to validity.

\par \textbf{Chapter~\ref{chap:conclusions} --- Conclusions} summarizes the contributions, answers the research questions based on the obtained results, and proposes future directions for research and development.

\chapter{Background and Related Work}
\label{chap:background}

\section{Classical Recommendation Systems}
\label{sec:bg-recsys}

\subsection{Collaborative Filtering and Co-Occurrence Analysis}
\label{subsec:bg-collab}

% Co-occurrence with Jaccard similarity, collaborative filtering fundamentals

\subsection{Content-Based Filtering}
\label{subsec:bg-content}

% Content-based similarity: category, tags, brand, price proximity

\subsection{Hybrid Recommendation Systems}
\label{subsec:bg-hybrid}

% Switching hybrid approach, cold-start strategies

\subsection{Trends with Temporal Decay}
\label{subsec:bg-trending}

% Exponential time decay for trending products

\section{Large Language Models in Recommendation Systems}
\label{sec:bg-llm-recsys}

\subsection{Roles of LLMs in Recommendation}
\label{subsec:bg-llm-roles}

% Knowledge enhancement, interaction, model enhancement (LLMEnhancedRecSys2024)

\subsection{Explainability of Recommendations via LLMs}
\label{subsec:bg-explainability}

% Natural language explanations, knowledge graphs (ExplainableRecKG2024)

\subsection{Two-Stage Architectures}
\label{subsec:bg-two-stage}

% Static candidate generation + LLM reranking/explanation

\section{Trust and Human-AI Interaction}
\label{sec:bg-trust}

\subsection{Proactive vs.\ Reactive Assistance}
\label{subsec:bg-proactive-reactive}

% CHI 2024 research, proactive vs reactive personalization

\subsection{Consumer Resistance to AI Chatbots}
\label{subsec:bg-resistance}

% Consumer distrust data, perceived intrusiveness

\subsection{Behavioral Signals as Intent Indicators}
\label{subsec:bg-behavior-signals}

% Cursor hover, dwell time, scroll depth as engagement proxies (CHI2011, SIGIR2020)

\section{Case Study: Amazon Rufus}
\label{sec:bg-rufus}

\subsection{System Architecture and Capabilities}
\label{subsec:bg-rufus-arch}

% Multi-model routing, RAG pipeline, streaming + hydration, COSMO knowledge graph

\subsection{Documented Accuracy and Trust Issues}
\label{subsec:bg-rufus-problems}

% 32% accuracy, 28% price hallucination, 83% self-serving, review fabrication

\subsection{Feature Evolution and Lessons Learned}
\label{subsec:bg-rufus-lessons}

% Reactive first, proactive added 10 months later; lessons for zalem

\section{Review Summarization and Product Comparison}
\label{sec:bg-review-comparison}

% Review summarization as decision support, AI-assisted product comparison

\chapter{Problem Definition and Requirements}
\label{chap:problem}

\section{Problem Statement}
\label{sec:prob-statement}

% E-commerce decision overload, limitations of existing AI assistants

\section{System Goals}
\label{sec:prob-goals}

% System goals and non-goals

\section{Functional Requirements}
\label{sec:prob-functional}

% Store browsing, recommendations, reactive advisor, review summarization, comparison

\section{Non-Functional Requirements}
\label{sec:prob-nonfunctional}

% Latency, cost, trust, explainability, privacy

\section{Evaluation Criteria}
\label{sec:prob-evaluation-criteria}

% Recommendation quality, system performance, user perception

\section{Constraints and Assumptions}
\label{sec:prob-constraints}

% Simulated store, seed data, limited user study, no real payments

\chapter{System Design}
\label{chap:design}

\section{Overall Architecture}
\label{sec:design-architecture}

% Three-layer architecture: client, Convex backend, external APIs

\section{Two-Stage Recommendation Pipeline}
\label{sec:design-pipeline}

\subsection{Stage 1: Candidate Generation}
\label{subsec:design-candidates}

% Co-occurrence (Jaccard), content similarity, trending, personalized hybrid
% <50ms target

\subsection{Stage 2: LLM Re-Ranking and Explanation}
\label{subsec:design-llm-reranking}

% Gemini Flash reranking, natural language message generation
% 300-500ms target

\section{Reactive-First Interaction Model}
\label{sec:design-reactive}

\subsection{Readiness Signals}
\label{subsec:design-readiness}

% Product dwell, deep scroll, review engagement, comparison behavior, cart deliberation

\subsection{Subtle Indicators and Question Chips}
\label{subsec:design-indicators}

% Breathing button, contextual chips, cooldown mechanism

\subsection{Privacy Boundary}
\label{subsec:design-privacy}

% Client-side aggregation, no raw events to backend

\section{Behavioral Tracking}
\label{sec:design-behavior}

% useDwellTime, useViewportTracking, useScrollDepth, useProductEngagement
% Event buffer with 5-second flush, composite interest score

\section{Review Summarization}
\label{sec:design-reviews}

% Batch processing with Flash-Lite, conflict surfacing, theme extraction

\section{AI-Assisted Comparison Mode}
\label{sec:design-comparison}

% 2-3 products, structured analysis, tradeoff identification

\section{Data Model}
\label{sec:design-data}

% Convex schema: products, categories, reviews, orders, cartItems, favorites,
% productCoOccurrences, behaviorSessions, suggestions, userPreferences

\section{Model Routing and Cost Strategy}
\label{sec:design-model-routing}

% Flash-Lite for batch/simple, Flash for conversational
% Context caching, token budget optimization

\section{Output Validation Strategy}
\label{sec:design-validation}

% Hydration over generation: LLM outputs productIds, client hydrates live data
% Never let LLM generate prices/specs

\chapter{Implementation}
\label{chap:implementation}

\section{Technology Stack}
\label{sec:impl-stack}

% Next.js 16, React 19, Convex, Clerk, Tailwind v4, Gemini Flash
% Turborepo monorepo with bun workspaces

\section{Project Structure}
\label{sec:impl-structure}

% apps/web, packages/backend, packages/ui, packages/env, packages/config

\section{E-Commerce Store}
\label{sec:impl-store}

\subsection{Application Pages}
\label{subsec:impl-pages}

% Home, category, product detail, search, cart, favorites, order history, checkout

\subsection{UI Components and Shared Library}
\label{subsec:impl-components}

% @zalem/ui with shadcn, ProductCard, ProductGrid, ProductFilters, etc.

\subsection{Interactions and UX Patterns}
\label{subsec:impl-ux}

% Optimistic updates, URL-based filters, skeleton loaders, responsive design

\section{Convex Data Layer}
\label{sec:impl-data}

\subsection{Schema and Indexing}
\label{subsec:impl-schema}

% Tables, indexes, denormalization strategy

\subsection{Seed Data}
\label{subsec:impl-seed}

% DummyJSON + faker, persona clusters, realistic order patterns

\section{Recommendation Engine}
\label{sec:impl-recommendations}

\subsection{Co-Occurrence with Jaccard Similarity}
\label{subsec:impl-jaccard}

% Implementation details, cron precomputation

\subsection{Content-Based Similarity}
\label{subsec:impl-content-sim}

% Weighted scoring formula

\subsection{Trending with Exponential Decay}
\label{subsec:impl-trending}

% Score formula, daily cron

\subsection{Personalized Hybrid Recommendations}
\label{subsec:impl-personalized}

% Switching hybrid logic, cold-start fallbacks

\subsection{Diversification with Maximal Marginal Relevance}
\label{subsec:impl-mmr}

% MMR for preventing same-brand clusters

\section{Behavioral Tracking}
\label{sec:impl-behavior}

% React hooks, event buffer, beacon flush, behaviorSessions

\section{Readiness Signals}
\label{sec:impl-readiness}

% Client-side evaluator, thresholds, cooldowns, question chips

\section{LLM Integration}
\label{sec:impl-llm}

\subsection{Agent Infrastructure}
\label{subsec:impl-agent}

% @convex-dev/agent, Gemini configuration, system prompt

\subsection{Advisor Sidebar Interface}
\label{subsec:impl-sidebar}

% Push-content layout, streaming, prompt-kit components

\subsection{Agent Tools}
\label{subsec:impl-tools}

% getProductDetails, searchProducts, getRecommendations, getCartContents, getReviewsSummary

\chapter{Evaluation}
\label{chap:evaluation}

\section{Evaluation Methodology}
\label{sec:eval-methodology}

% Three-level evaluation: recommendation quality, system performance, user perception

\section{Recommendation System Evaluation}
\label{sec:eval-recsys}

\subsection{Offline Evaluation Protocol}
\label{subsec:eval-offline-protocol}

% Train/test split, holdout setup

\subsection{Metrics}
\label{subsec:eval-metrics}

% Hit Rate@5, NDCG@10, catalog coverage, diversity, popularity bias

\subsection{Results}
\label{subsec:eval-recsys-results}

% Tables and charts

\section{System Performance Evaluation}
\label{sec:eval-system}

\subsection{Latency}
\label{subsec:eval-latency}

% p50, p95 for recommendation queries, LLM responses, comparison generation

\subsection{Cost and Token Usage}
\label{subsec:eval-cost}

% Token usage per request, estimated cost per session, cache hit impact

\subsection{Results}
\label{subsec:eval-system-results}

% Tables and charts

\section{User Evaluation}
\label{sec:eval-user}

\subsection{Participant Profile}
\label{subsec:eval-participants}

% Number, demographics, technical background

\subsection{Study Design}
\label{subsec:eval-study-design}

% Tasks, conditions (baseline vs reactive-first), within-subjects

\subsection{Measurement Instruments}
\label{subsec:eval-instruments}

% Likert scales, SUS, trust/adoption questions

\subsection{Results}
\label{subsec:eval-user-results}

% Perceived usefulness, trust, intrusiveness, purchase confidence
% Review summarization usefulness, comparison usefulness

\section{Discussion}
\label{sec:eval-discussion}

\subsection{Interpretation of Results}
\label{subsec:eval-interpretation}

% Did reactive-first help? Where did explanations help?
% Review summaries, comparison mode effectiveness

\subsection{Limitations and Threats to Validity}
\label{subsec:eval-limitations}

% Limited participants, synthetic data, bounded domain, novelty bias

\chapter{Conclusions}
\label{chap:conclusions}

\section{Summary of Contributions}
\label{sec:concl-summary}

% Recap: hybrid architecture, reactive-first model, review summarization,
% comparison mode, behavior-driven readiness, evaluation

\section{Answers to Research Questions}
\label{sec:concl-rq-answers}

% RQ1-RQ4 answered based on evaluation results

\section{Practical Implications}
\label{sec:concl-implications}

% Design guidelines for non-intrusive AI shopping assistants

\section{Future Work}
\label{sec:concl-future}

% Larger-scale evaluation, broader product domains, real user data,
% advanced behavior signals (eye tracking), A/B testing


\bibliography{references}

\end{document}
